\section{Metodología}
\label{sec:metodologia}

\subsection{Los modelos y qué se hizo con los pesos}
\label{sec:pesos}

Esta subsección es deliberadamente explícita, porque el tratamiento de los
pesos pre-entrenados es donde un experimento de \textit{fine-tuning} se
invalida con más facilidad y sin dar ningún error.

\paragraph{Origen.} Los dos modelos se construyen a partir de los
\textit{checkpoints} públicos de HuggingFace, que se enumeran en la
Tabla~\ref{tab:modelos}. DistilBERT es una destilación del primero de ellos:
conserva una de cada dos capas, elimina el \textit{pooler} y los
\textit{embeddings} de segmento, y mantiene la dimensión oculta (768) y el
vocabulario. Los pesos se descargan del Hub en el momento de construir el
modelo; \textbf{no hay ningún peso almacenado en el repositorio}, ni
versionado, ni modificado, ni redistribuido. El repositorio contiene
únicamente código, métricas en JSON y figuras.

\begin{table}[h]
\centering
\small
\begin{tabular}{llcr}
\toprule
\textbf{Nombre corto} & \textbf{\textit{Checkpoint} en el Hub} &
\textbf{Capas} & \textbf{Parámetros} \\
\midrule
BERT-base       & \texttt{bert-base-uncased}       & 12 & \num{109483778} \\
DistilBERT-base & \texttt{distilbert-base-uncased} &  6 & \num{66955010}  \\
\bottomrule
\end{tabular}
\caption{Los dos \textit{checkpoints} de partida. El número de parámetros
corresponde a una cabeza de dos clases; con las cuatro clases de AG News
sube en \num{1538} en ambos modelos ($768 \times 2 + 2$ pesos de salida
adicionales).}
\label{tab:modelos}
\end{table}

\paragraph{Qué se inicializa al azar.} Al construir un clasificador, el
\textit{encoder} llega pre-entrenado pero la cabeza de clasificación se
inicializa aleatoriamente: es la única parte que no ha visto nunca datos. El
aviso \textit{``newly initialized''} que imprime \texttt{transformers} es
esperado y correcto; su ausencia sería la señal de alarma. La semilla 42 fija
esa inicialización, de modo que la cabeza de BERT y la de DistilBERT parten
del mismo estado aleatorio.

\paragraph{Qué se actualiza durante el entrenamiento.} En las corridas base
se actualizan todos los parámetros: \textit{encoder} y cabeza. En las
configuraciones congeladas del estudio de ablación se pone
\texttt{requires\_grad = False} sobre el \textit{encoder} y, además, esos
parámetros \textbf{no se pasan al optimizador}. Los dos pasos son necesarios:
con solo el primero, AdamW seguiría aplicando \textit{weight decay} a pesos
que no reciben gradiente, y el modelo se degradaría de forma silenciosa.

\paragraph{El tokenizador se deriva siempre del modelo.} Cada
\textit{checkpoint} trae su propio vocabulario. Usar el tokenizador de otro
modelo rompe la correspondencia entre \textit{token} e \textit{embedding} y
el modelo predice ruido sin lanzar ninguna excepción. Por eso
\texttt{build\_tokenizer()} nunca acepta un vocabulario suelto: lo deriva del
identificador del modelo.

\paragraph{Qué se conserva y qué no.} Durante el entrenamiento se mantiene en
memoria una copia del \textit{state dict} de la época con mejor F1 de
validación, y es esa copia la que se carga antes de evaluar en test. Al
terminar, el modelo ajustado \textbf{se descarta}: no se escribe en disco ni
se versiona. Todo lo que sobrevive a una corrida son sus métricas. Es una
decisión consciente --- los \textit{checkpoints} ajustados pesan cientos de
megabytes cada uno y son reproducibles desde el código y la semilla --- y
tiene la consecuencia práctica de que cargar de nuevo cualquier modelo del
proyecto parte siempre de los pesos originales del Hub, sin estados
intermedios que puedan haberse corrompido.

\subsection{La cabeza de clasificación del estudio de ablación}

Las clases \texttt{AutoModelForSequenceClassification} de HuggingFace traen
una cabeza fija, y --- esto es lo que obliga a no usarlas aquí --- una cabeza
\textbf{distinta en cada arquitectura}: BERT pasa el \textit{token}
\texttt{[CLS]} por un \textit{pooler} con \texttt{tanh}; DistilBERT usa un
\texttt{pre\_classifier} con \texttt{ReLU}. Comparar cabezas montadas sobre
dos preprocesos distintos mezclaría dos variables en el mismo experimento.

Por eso el estudio de ablación usa una envoltura propia
(\texttt{ClasificadorConCabeza} en \texttt{src/models.py}): \textit{encoder}
desnudo (\texttt{AutoModel}) más un MLP construido por nosotros sobre el
estado oculto del \textit{token} \texttt{[CLS]}, idéntico para las dos
arquitecturas. Así la única diferencia entre dos corridas es lo que el
experimento cambia a propósito.

\subsection{Protocolo de entrenamiento}

El bucle de entrenamiento está escrito a mano, sin el \texttt{Trainer} de
HuggingFace, para tener control explícito sobre el \textit{scheduler}, el
recorte de gradientes y el registro por iteración. Los hiperparámetros son
idénticos para los dos modelos y las tres tareas.

\begin{table}[h]
\centering
\small
\begin{tabular}{ll}
\toprule
Optimizador        & AdamW, \texttt{lr} $= 2\times10^{-5}$, \texttt{weight\_decay} $= 0{,}01$ \\
\textit{Weight decay} & No se aplica a \textit{bias} ni a \texttt{LayerNorm} \\
\textit{Scheduler} & \textit{Warmup} lineal \SI{10}{\percent} + decaimiento lineal \\
Épocas / \textit{batch} & 2 / 32 \\
Precisión          & fp16 (\textit{mixed precision}) \\
Recorte            & Norma máxima de gradiente $= 1{,}0$ \\
Selección          & \textit{Checkpoint} con mejor F1 de \textbf{validación} \\
Semilla            & 42 en Python, NumPy y PyTorch \\
\bottomrule
\end{tabular}
\caption{Protocolo común. El \textit{learning rate} está dentro del rango
recomendado en el artículo original de BERT ($5\times10^{-5}$,
$3\times10^{-5}$ o $2\times10^{-5}$).}
\end{table}

Dos detalles del bucle que cambian los resultados si se omiten. El primero:
el \textit{weight decay} no se aplica a los términos de \textit{bias} ni a los
de \texttt{LayerNorm}, porque son parámetros de escala y desplazamiento y
penalizarlos degrada el modelo. El segundo: con precisión mixta hay que
des-escalar los gradientes \emph{antes} de recortarlos
(\texttt{scaler.unscale\_()} antes de \texttt{clip\_grad\_norm\_()}), o el
umbral se aplicaría sobre gradientes escalados y no significaría lo que se
cree.

\paragraph{El test se mira una sola vez.} La configuración ganadora del
estudio de ablación se elige por F1 de validación. El conjunto de test se
evalúa al final, con el \textit{checkpoint} ya seleccionado, y no participa en
ninguna decisión. Elegir sobre el conjunto que después se reporta es la forma
más común de producir números optimistas.

\subsection{Métricas}

\textit{Precision}, \textit{recall} y F1 se reportan como promedio macro: se
promedia la métrica de cada clase dándoles el mismo peso. Es lo adecuado con
clases desbalanceadas (SST-2 está en 56/44) y, además, el promedio micro en
clasificación multiclase de etiqueta única colapsa a la \textit{accuracy} y no
aportaría información.

Para la eficiencia se miden parámetros, tamaño en fp32, latencia de inferencia
y pico de memoria de GPU. La medición de latencia toma dos precauciones que
cambian el resultado por completo: un \textit{warm-up} de varias pasadas
--- las primeras incluyen la carga de \textit{kernels} CUDA y la reserva de
memoria, y medirlas infla la cifra varias veces --- y una llamada a
\texttt{torch.cuda.synchronize()} en cada iteración, porque las llamadas a
CUDA son asíncronas y sin sincronizar se mediría el tiempo de \emph{encolar}
la operación, no el de ejecutarla. Para la memoria se resetea el contador de
pico antes de medir, o se arrastraría el pico del entrenamiento previo.
